\documentclass[10pt,a4paper]{article}
\usepackage[utf8]{inputenc}
\usepackage{amsmath,amssymb,amsfonts}
\usepackage{graphicx}
\usepackage{hyperref}
\usepackage{bm}
\usepackage{booktabs}
\usepackage{algorithm}
\usepackage{algorithmicx}
\usepackage{algpseudocode}

\title{Gravit Grover: Distributed Epistemic Search with Gower Geometry and Amplification Dynamics}
\author{GravitOpenNetwork Team}
\date{\today}

\begin{document}

\maketitle

\begin{abstract}
We introduce Gravit Grover, a distributed epistemic search protocol for converging on globally consistent hypotheses in decentralized networks. The protocol combines Gower-based similarity geometry with multiplicative weights updates and Grover-like amplification to achieve consensus without central coordination. We formalize the four-layer architecture: (1) Gower Similarity Layer for geometric representation of hypothesis space, (2) Probabilistic Layer for Bayesian-style belief updates, (3) Grover Amplification Layer for exponential enhancement of consistent hypotheses, and (4) Gossip Consensus Layer for distributed synchronization. We prove convergence guarantees through spectral gap analysis and demonstrate the protocol's effectiveness through empirical experiments. Gravit Grover provides a mathematically grounded framework for decentralized AI reasoning, enabling verifiable consensus formation across heterogeneous validator nodes.
\end{abstract}

\section{Introduction}

The rapid advancement of large language models (LLMs) has created both opportunities and challenges for AI-assisted decision making. While LLMs demonstrate remarkable capabilities in generating plausible hypotheses, their stochastic nature and susceptibility to hallucinations necessitate validation mechanisms. Traditional approaches rely on centralized evaluation or simple voting, which introduce single points of failure and are vulnerable to manipulation.

\subsection{Motivation}

We identify three key requirements for robust hypothesis evaluation in distributed systems:

\begin{enumerate}
    \item \textbf{Verifiability:} The reasoning process should be transparent and reproducible.
    \item \textbf{Decentralization:} No single node should control the evaluation outcome.
    \item \textbf{Convergence:} The system should mathematically converge to optimal hypotheses.
\end{enumerate}

\subsection{Contributions}

Our main contributions are:

\begin{enumerate}
    \item A novel protocol combining Gower geometry with multiplicative weights and amplification.
    \item Formal convergence guarantees through spectral graph theory.
    \item A practical implementation with support for LLM-based validators.
    \item Empirical validation on synthetic and real-world tasks.
\end{enumerate}

\section{Mathematical Framework}

\subsection{Gower Similarity Layer}

We define the hypothesis space $\mathcal{H}$ with features $x_i \in \mathbb{R}^d$ for each hypothesis $h_i$. The Gower similarity between hypotheses $h_i$ and $h_j$ is:

\begin{equation}
S(h_i, h_j) = \frac{\sum_{k} w_k \cdot s_k(h_i, h_j)}{\sum_{k} w_k}
\end{equation}

where $s_k$ is the similarity on feature $k$ and $w_k$ is the feature weight.

For numeric features:
\begin{equation}
s_k(h_i, h_j) = 1 - \frac{|x_{ik} - x_{jk}|}{R_k}
\end{equation}

For categorical features:
\begin{equation}
s_k(h_i, h_j) = \mathbb{I}[x_{ik} = x_{jk}]
\end{equation}

\subsection{Multiplicative Weights Update}

Each node $v$ maintains a probability distribution $p_v$ over hypotheses. The update rule is:

\begin{equation}
p_v^{t+1}(h) = \frac{p_v^t(h) \cdot \exp(\eta \cdot \text{score}_v(h))}{Z_t}
\end{equation}

where $\eta$ is the learning rate and $Z_t$ is the normalization constant.

\subsection{Grover Amplification}

The amplified distribution is computed as:

\begin{equation}
p^{t+1}(h) \propto p^t(h) \cdot \left(\frac{\text{global\_score}(h)}{\text{avg\_score}}\right)^\gamma
\end{equation}

where $\gamma > 1$ controls the amplification strength.

\subsection{Gossip Consensus}

Consensus is achieved through iterative mixing:

\begin{equation}
x_v^{t+1} = (1 - \epsilon)x_v^t + \frac{\epsilon}{|N(v)|} \sum_{u \in N(v)} x_u^t
\end{equation}

\section{Convergence Analysis}

\section{Implementation}

\section{Experiments}

\section{Conclusion}

\bibliography{references}
\bibliographystyle{plain}

\end{document}
